\documentclass[11pt,a4paper]{article}
\newcommand{\gradelabel}{Grade 11}
\newcommand{\leconhdr}{Capteurs numériques et analogiques}
% Préambule commun aux plans de cours de robotique (Grades 7 à 12).
% Chaque plan de cours commence par \documentclass[11pt,a4paper]{article}, définit
% \gradelabel et \leconhdr (en-têtes de page), puis % Préambule commun aux plans de cours de robotique (Grades 7 à 12).
% Chaque plan de cours commence par \documentclass[11pt,a4paper]{article}, définit
% \gradelabel et \leconhdr (en-têtes de page), puis % Préambule commun aux plans de cours de robotique (Grades 7 à 12).
% Chaque plan de cours commence par \documentclass[11pt,a4paper]{article}, définit
% \gradelabel et \leconhdr (en-têtes de page), puis \input{../../preambule_robotique.tex}.
\usepackage[T1]{fontenc}
\usepackage[utf8]{inputenc}
\usepackage{lmodern}
\usepackage[french]{babel}
\usepackage[margin=2cm,headheight=14pt]{geometry}
\usepackage[most]{tcolorbox}
\usepackage{xcolor,tabularx,booktabs,array,enumitem,fancyhdr,amssymb,tikz}
\usetikzlibrary{arrows.meta,positioning}

\pagestyle{fancy}
\fancyhf{}
\lhead{\small Robotique -- \gradelabel{} -- Novembre 2026}
\rhead{\small \leconhdr}
\cfoot{\thepage}
\setlist{itemsep=1pt,topsep=2pt,leftmargin=*}
\renewcommand{\arraystretch}{1.35}
\newcolumntype{Y}{>{\raggedright\arraybackslash}X}
\newcommand{\cb}{$\square$\,}
\newcommand{\fl}{$\rightarrow$}
\newcommand{\lig}{\noindent\rule{\linewidth}{0.3pt}\\[5mm]}

\newtcolorbox{etape}[3]{enhanced,breakable,colback=#1!6,colframe=#1!70!black,
  coltitle=white,fonttitle=\bfseries,title={#2\hfill #3},boxrule=0.8pt,arc=2mm,
  before skip=8pt,after skip=8pt}
\newtcolorbox{astuce}[1][Différenciation et besoins particuliers]{breakable,
  colback=yellow!12,colframe=yellow!50!black,fonttitle=\bfseries\small,
  title={#1},boxrule=0.5pt,arc=1mm,before skip=6pt,after skip=6pt}
\newtcolorbox{code}[1][Code]{breakable,colback=black!4,colframe=black!60,
  fonttitle=\bfseries\small,title={#1},boxrule=0.5pt,arc=1mm,fontupper=\ttfamily\small,
  before skip=6pt,after skip=6pt}
% Blocs de code coloré (écrire le code sans accents : comme dans l'éditeur)
\newtcblisting{arduino}[1][Code Arduino]{listing only,breakable,colback=black!4,
  colframe=black!60,fonttitle=\bfseries\small,title={#1},boxrule=0.5pt,arc=1mm,
  before skip=6pt,after skip=6pt,
  listing options={language=C++,basicstyle=\ttfamily\small,columns=fullflexible,
  keepspaces=true,showstringspaces=false,commentstyle=\color{green!40!black}\itshape,
  keywordstyle=\color{blue!70!black}\bfseries}}
\newtcblisting{python}[1][Code Python]{listing only,breakable,colback=black!4,
  colframe=black!60,fonttitle=\bfseries\small,title={#1},boxrule=0.5pt,arc=1mm,
  before skip=6pt,after skip=6pt,
  listing options={language=Python,basicstyle=\ttfamily\small,columns=fullflexible,
  keepspaces=true,showstringspaces=false,commentstyle=\color{green!40!black}\itshape,
  keywordstyle=\color{blue!70!black}\bfseries}}
\newcommand{\entete}[2]{%
  \begin{tcolorbox}[colback=#1!10,colframe=#1!70!black,boxrule=0.8pt,arc=2mm]
  {\Large\bfseries #2}\hfill Nom : \rule{4cm}{0.3pt}\quad Date : \rule{2.2cm}{0.3pt}
  \end{tcolorbox}}

% Bandeau de titre : #1 titre, #2 semaine
\newcommand{\titreblock}[2]{%
\begin{tcolorbox}[colback=teal!10,colframe=teal!70!black,arc=2mm,boxrule=1pt]
\centering
{\LARGE\bfseries #1}\\[2mm]
{\large Robotique -- \gradelabel{} -- Plan de cours de 2 heures}\\[1mm]
Enseignant : Francis Ki \quad|\quad #2
\end{tcolorbox}}

% Tableau du minutage : 6 contenus (accroche, mini-leçon, guidée, autonome, mise en commun, exit ticket)
\newcommand{\minutage}[6]{%
\begin{tabularx}{\linewidth}{>{\bfseries}p{2.2cm}p{2.6cm}Y p{3cm}}
\toprule
Minutes & Étape & Contenu & Organisation\\
\midrule
0 -- 10 & Accroche & #1 & Classe entière\\
10 -- 30 & Mini-leçon & #2 & Classe entière\\
30 -- 60 & Pratique guidée & #3 & Collectif puis binômes\\
60 -- 95 & Pratique autonome & #4 & Individuel sur poste\\
95 -- 110 & Mise en commun & #5 & Classe entière\\
110 -- 120 & Exit ticket & #6 & Individuel sur poste\\
\bottomrule
\end{tabularx}}

% Étayage progressif : #1 aide 1, #2 aide 2, #3 aide 3
\newcommand{\aides}[3]{%
\textbf{Étayage progressif de l'aide.}
\begin{center}
\begin{tabularx}{\linewidth}{>{\bfseries}p{2.4cm}Y p{4.4cm}}
\toprule
Niveau d'aide & Contenu & Qui intervient\\
\midrule
Aide 1 & #1 & Le binôme seul\\
Aide 2 & #2 & Le binôme avec la carte d'aide\\
Aide 3 & #3 & L'enseignant\\
\bottomrule
\end{tabularx}
\end{center}}

% Les trois groupes de la pratique autonome : #1 tâche A, #2 tâche B, #3 tâche C
\newcommand{\groupes}[3]{%
\textbf{Consigne commune\,:} chaque élève ouvre le fichier de son groupe, le renomme avec son prénom et travaille seul sur son poste. L'enseignant passe d'abord auprès du groupe A pendant 10 min.

\medskip
\begin{tabularx}{\linewidth}{>{\bfseries}p{3cm}Y p{2.2cm}}
\toprule
Groupe & Tâche sur poste & Annexe\\
\midrule
A : en dessous du niveau & #1 & Annexe C\\
B : au niveau & #2 & Annexe D\\
C : au-dessus du niveau & #3 & Annexe E\\
\bottomrule
\end{tabularx}}

% Mise en commun : #1 groupe A, #2 groupe B, #3 groupe C, #4 conclusion
\newcommand{\miseencommun}[4]{%
\begin{etape}{teal}{5. Mise en commun}{95 -- 110 min}
\begin{enumerate}
 \item (2 min) Retour au calme, écrans mis en veille.
 \item (9 min) Trois passages courts de 3 min, avec valorisation :
 \begin{itemize}
  \item \textbf{Groupe A} : #1
  \item \textbf{Groupe B} : #2
  \item \textbf{Groupe C} : #3
 \end{itemize}
 \item (4 min) Conclusion : #4
\end{enumerate}
\end{etape}}

% Clôture : #1 question 1, #2 question 2
\newcommand{\cloture}[2]{%
\begin{etape}{gray}{6. Clôture : exit ticket}{110 -- 120 min}
Formulaire en ligne ou papier (annexe G) :
\begin{enumerate}
 \item #1
 \item #2
 \item Auto-évaluation : \emph{Je maîtrise / Je progresse / J'ai besoin d'aide.}
\end{enumerate}
Lire les réponses avant la séance suivante pour ajuster les groupes de besoin.
\end{etape}}

% Grille d'observation : #1 #2 #3 intitulés de colonnes
\newcommand{\grille}[3]{%
\subsection*{Grille d'observation (enseignant)}
\begin{tabularx}{\linewidth}{Y p{2.4cm} p{2.4cm} p{2.4cm} Y}
\toprule
Élève & #1 & #2 & #3 & Groupe suivant\\
\midrule
 & & & & \\ & & & & \\ & & & & \\ & & & & \\ & & & & \\
\bottomrule
\end{tabularx}}

% Signature de l'exit ticket
\newcommand{\autoeval}{Je me sens : \cb Je maîtrise \quad \cb Je progresse \quad \cb J'ai besoin d'aide}
.
\usepackage[T1]{fontenc}
\usepackage[utf8]{inputenc}
\usepackage{lmodern}
\usepackage[french]{babel}
\usepackage[margin=2cm,headheight=14pt]{geometry}
\usepackage[most]{tcolorbox}
\usepackage{xcolor,tabularx,booktabs,array,enumitem,fancyhdr,amssymb,tikz}
\usetikzlibrary{arrows.meta,positioning}

\pagestyle{fancy}
\fancyhf{}
\lhead{\small Robotique -- \gradelabel{} -- Novembre 2026}
\rhead{\small \leconhdr}
\cfoot{\thepage}
\setlist{itemsep=1pt,topsep=2pt,leftmargin=*}
\renewcommand{\arraystretch}{1.35}
\newcolumntype{Y}{>{\raggedright\arraybackslash}X}
\newcommand{\cb}{$\square$\,}
\newcommand{\fl}{$\rightarrow$}
\newcommand{\lig}{\noindent\rule{\linewidth}{0.3pt}\\[5mm]}

\newtcolorbox{etape}[3]{enhanced,breakable,colback=#1!6,colframe=#1!70!black,
  coltitle=white,fonttitle=\bfseries,title={#2\hfill #3},boxrule=0.8pt,arc=2mm,
  before skip=8pt,after skip=8pt}
\newtcolorbox{astuce}[1][Différenciation et besoins particuliers]{breakable,
  colback=yellow!12,colframe=yellow!50!black,fonttitle=\bfseries\small,
  title={#1},boxrule=0.5pt,arc=1mm,before skip=6pt,after skip=6pt}
\newtcolorbox{code}[1][Code]{breakable,colback=black!4,colframe=black!60,
  fonttitle=\bfseries\small,title={#1},boxrule=0.5pt,arc=1mm,fontupper=\ttfamily\small,
  before skip=6pt,after skip=6pt}
% Blocs de code coloré (écrire le code sans accents : comme dans l'éditeur)
\newtcblisting{arduino}[1][Code Arduino]{listing only,breakable,colback=black!4,
  colframe=black!60,fonttitle=\bfseries\small,title={#1},boxrule=0.5pt,arc=1mm,
  before skip=6pt,after skip=6pt,
  listing options={language=C++,basicstyle=\ttfamily\small,columns=fullflexible,
  keepspaces=true,showstringspaces=false,commentstyle=\color{green!40!black}\itshape,
  keywordstyle=\color{blue!70!black}\bfseries}}
\newtcblisting{python}[1][Code Python]{listing only,breakable,colback=black!4,
  colframe=black!60,fonttitle=\bfseries\small,title={#1},boxrule=0.5pt,arc=1mm,
  before skip=6pt,after skip=6pt,
  listing options={language=Python,basicstyle=\ttfamily\small,columns=fullflexible,
  keepspaces=true,showstringspaces=false,commentstyle=\color{green!40!black}\itshape,
  keywordstyle=\color{blue!70!black}\bfseries}}
\newcommand{\entete}[2]{%
  \begin{tcolorbox}[colback=#1!10,colframe=#1!70!black,boxrule=0.8pt,arc=2mm]
  {\Large\bfseries #2}\hfill Nom : \rule{4cm}{0.3pt}\quad Date : \rule{2.2cm}{0.3pt}
  \end{tcolorbox}}

% Bandeau de titre : #1 titre, #2 semaine
\newcommand{\titreblock}[2]{%
\begin{tcolorbox}[colback=teal!10,colframe=teal!70!black,arc=2mm,boxrule=1pt]
\centering
{\LARGE\bfseries #1}\\[2mm]
{\large Robotique -- \gradelabel{} -- Plan de cours de 2 heures}\\[1mm]
Enseignant : Francis Ki \quad|\quad #2
\end{tcolorbox}}

% Tableau du minutage : 6 contenus (accroche, mini-leçon, guidée, autonome, mise en commun, exit ticket)
\newcommand{\minutage}[6]{%
\begin{tabularx}{\linewidth}{>{\bfseries}p{2.2cm}p{2.6cm}Y p{3cm}}
\toprule
Minutes & Étape & Contenu & Organisation\\
\midrule
0 -- 10 & Accroche & #1 & Classe entière\\
10 -- 30 & Mini-leçon & #2 & Classe entière\\
30 -- 60 & Pratique guidée & #3 & Collectif puis binômes\\
60 -- 95 & Pratique autonome & #4 & Individuel sur poste\\
95 -- 110 & Mise en commun & #5 & Classe entière\\
110 -- 120 & Exit ticket & #6 & Individuel sur poste\\
\bottomrule
\end{tabularx}}

% Étayage progressif : #1 aide 1, #2 aide 2, #3 aide 3
\newcommand{\aides}[3]{%
\textbf{Étayage progressif de l'aide.}
\begin{center}
\begin{tabularx}{\linewidth}{>{\bfseries}p{2.4cm}Y p{4.4cm}}
\toprule
Niveau d'aide & Contenu & Qui intervient\\
\midrule
Aide 1 & #1 & Le binôme seul\\
Aide 2 & #2 & Le binôme avec la carte d'aide\\
Aide 3 & #3 & L'enseignant\\
\bottomrule
\end{tabularx}
\end{center}}

% Les trois groupes de la pratique autonome : #1 tâche A, #2 tâche B, #3 tâche C
\newcommand{\groupes}[3]{%
\textbf{Consigne commune\,:} chaque élève ouvre le fichier de son groupe, le renomme avec son prénom et travaille seul sur son poste. L'enseignant passe d'abord auprès du groupe A pendant 10 min.

\medskip
\begin{tabularx}{\linewidth}{>{\bfseries}p{3cm}Y p{2.2cm}}
\toprule
Groupe & Tâche sur poste & Annexe\\
\midrule
A : en dessous du niveau & #1 & Annexe C\\
B : au niveau & #2 & Annexe D\\
C : au-dessus du niveau & #3 & Annexe E\\
\bottomrule
\end{tabularx}}

% Mise en commun : #1 groupe A, #2 groupe B, #3 groupe C, #4 conclusion
\newcommand{\miseencommun}[4]{%
\begin{etape}{teal}{5. Mise en commun}{95 -- 110 min}
\begin{enumerate}
 \item (2 min) Retour au calme, écrans mis en veille.
 \item (9 min) Trois passages courts de 3 min, avec valorisation :
 \begin{itemize}
  \item \textbf{Groupe A} : #1
  \item \textbf{Groupe B} : #2
  \item \textbf{Groupe C} : #3
 \end{itemize}
 \item (4 min) Conclusion : #4
\end{enumerate}
\end{etape}}

% Clôture : #1 question 1, #2 question 2
\newcommand{\cloture}[2]{%
\begin{etape}{gray}{6. Clôture : exit ticket}{110 -- 120 min}
Formulaire en ligne ou papier (annexe G) :
\begin{enumerate}
 \item #1
 \item #2
 \item Auto-évaluation : \emph{Je maîtrise / Je progresse / J'ai besoin d'aide.}
\end{enumerate}
Lire les réponses avant la séance suivante pour ajuster les groupes de besoin.
\end{etape}}

% Grille d'observation : #1 #2 #3 intitulés de colonnes
\newcommand{\grille}[3]{%
\subsection*{Grille d'observation (enseignant)}
\begin{tabularx}{\linewidth}{Y p{2.4cm} p{2.4cm} p{2.4cm} Y}
\toprule
Élève & #1 & #2 & #3 & Groupe suivant\\
\midrule
 & & & & \\ & & & & \\ & & & & \\ & & & & \\ & & & & \\
\bottomrule
\end{tabularx}}

% Signature de l'exit ticket
\newcommand{\autoeval}{Je me sens : \cb Je maîtrise \quad \cb Je progresse \quad \cb J'ai besoin d'aide}
.
\usepackage[T1]{fontenc}
\usepackage[utf8]{inputenc}
\usepackage{lmodern}
\usepackage[french]{babel}
\usepackage[margin=2cm,headheight=14pt]{geometry}
\usepackage[most]{tcolorbox}
\usepackage{xcolor,tabularx,booktabs,array,enumitem,fancyhdr,amssymb,tikz}
\usetikzlibrary{arrows.meta,positioning}

\pagestyle{fancy}
\fancyhf{}
\lhead{\small Robotique -- \gradelabel{} -- Novembre 2026}
\rhead{\small \leconhdr}
\cfoot{\thepage}
\setlist{itemsep=1pt,topsep=2pt,leftmargin=*}
\renewcommand{\arraystretch}{1.35}
\newcolumntype{Y}{>{\raggedright\arraybackslash}X}
\newcommand{\cb}{$\square$\,}
\newcommand{\fl}{$\rightarrow$}
\newcommand{\lig}{\noindent\rule{\linewidth}{0.3pt}\\[5mm]}

\newtcolorbox{etape}[3]{enhanced,breakable,colback=#1!6,colframe=#1!70!black,
  coltitle=white,fonttitle=\bfseries,title={#2\hfill #3},boxrule=0.8pt,arc=2mm,
  before skip=8pt,after skip=8pt}
\newtcolorbox{astuce}[1][Différenciation et besoins particuliers]{breakable,
  colback=yellow!12,colframe=yellow!50!black,fonttitle=\bfseries\small,
  title={#1},boxrule=0.5pt,arc=1mm,before skip=6pt,after skip=6pt}
\newtcolorbox{code}[1][Code]{breakable,colback=black!4,colframe=black!60,
  fonttitle=\bfseries\small,title={#1},boxrule=0.5pt,arc=1mm,fontupper=\ttfamily\small,
  before skip=6pt,after skip=6pt}
% Blocs de code coloré (écrire le code sans accents : comme dans l'éditeur)
\newtcblisting{arduino}[1][Code Arduino]{listing only,breakable,colback=black!4,
  colframe=black!60,fonttitle=\bfseries\small,title={#1},boxrule=0.5pt,arc=1mm,
  before skip=6pt,after skip=6pt,
  listing options={language=C++,basicstyle=\ttfamily\small,columns=fullflexible,
  keepspaces=true,showstringspaces=false,commentstyle=\color{green!40!black}\itshape,
  keywordstyle=\color{blue!70!black}\bfseries}}
\newtcblisting{python}[1][Code Python]{listing only,breakable,colback=black!4,
  colframe=black!60,fonttitle=\bfseries\small,title={#1},boxrule=0.5pt,arc=1mm,
  before skip=6pt,after skip=6pt,
  listing options={language=Python,basicstyle=\ttfamily\small,columns=fullflexible,
  keepspaces=true,showstringspaces=false,commentstyle=\color{green!40!black}\itshape,
  keywordstyle=\color{blue!70!black}\bfseries}}
\newcommand{\entete}[2]{%
  \begin{tcolorbox}[colback=#1!10,colframe=#1!70!black,boxrule=0.8pt,arc=2mm]
  {\Large\bfseries #2}\hfill Nom : \rule{4cm}{0.3pt}\quad Date : \rule{2.2cm}{0.3pt}
  \end{tcolorbox}}

% Bandeau de titre : #1 titre, #2 semaine
\newcommand{\titreblock}[2]{%
\begin{tcolorbox}[colback=teal!10,colframe=teal!70!black,arc=2mm,boxrule=1pt]
\centering
{\LARGE\bfseries #1}\\[2mm]
{\large Robotique -- \gradelabel{} -- Plan de cours de 2 heures}\\[1mm]
Enseignant : Francis Ki \quad|\quad #2
\end{tcolorbox}}

% Tableau du minutage : 6 contenus (accroche, mini-leçon, guidée, autonome, mise en commun, exit ticket)
\newcommand{\minutage}[6]{%
\begin{tabularx}{\linewidth}{>{\bfseries}p{2.2cm}p{2.6cm}Y p{3cm}}
\toprule
Minutes & Étape & Contenu & Organisation\\
\midrule
0 -- 10 & Accroche & #1 & Classe entière\\
10 -- 30 & Mini-leçon & #2 & Classe entière\\
30 -- 60 & Pratique guidée & #3 & Collectif puis binômes\\
60 -- 95 & Pratique autonome & #4 & Individuel sur poste\\
95 -- 110 & Mise en commun & #5 & Classe entière\\
110 -- 120 & Exit ticket & #6 & Individuel sur poste\\
\bottomrule
\end{tabularx}}

% Étayage progressif : #1 aide 1, #2 aide 2, #3 aide 3
\newcommand{\aides}[3]{%
\textbf{Étayage progressif de l'aide.}
\begin{center}
\begin{tabularx}{\linewidth}{>{\bfseries}p{2.4cm}Y p{4.4cm}}
\toprule
Niveau d'aide & Contenu & Qui intervient\\
\midrule
Aide 1 & #1 & Le binôme seul\\
Aide 2 & #2 & Le binôme avec la carte d'aide\\
Aide 3 & #3 & L'enseignant\\
\bottomrule
\end{tabularx}
\end{center}}

% Les trois groupes de la pratique autonome : #1 tâche A, #2 tâche B, #3 tâche C
\newcommand{\groupes}[3]{%
\textbf{Consigne commune\,:} chaque élève ouvre le fichier de son groupe, le renomme avec son prénom et travaille seul sur son poste. L'enseignant passe d'abord auprès du groupe A pendant 10 min.

\medskip
\begin{tabularx}{\linewidth}{>{\bfseries}p{3cm}Y p{2.2cm}}
\toprule
Groupe & Tâche sur poste & Annexe\\
\midrule
A : en dessous du niveau & #1 & Annexe C\\
B : au niveau & #2 & Annexe D\\
C : au-dessus du niveau & #3 & Annexe E\\
\bottomrule
\end{tabularx}}

% Mise en commun : #1 groupe A, #2 groupe B, #3 groupe C, #4 conclusion
\newcommand{\miseencommun}[4]{%
\begin{etape}{teal}{5. Mise en commun}{95 -- 110 min}
\begin{enumerate}
 \item (2 min) Retour au calme, écrans mis en veille.
 \item (9 min) Trois passages courts de 3 min, avec valorisation :
 \begin{itemize}
  \item \textbf{Groupe A} : #1
  \item \textbf{Groupe B} : #2
  \item \textbf{Groupe C} : #3
 \end{itemize}
 \item (4 min) Conclusion : #4
\end{enumerate}
\end{etape}}

% Clôture : #1 question 1, #2 question 2
\newcommand{\cloture}[2]{%
\begin{etape}{gray}{6. Clôture : exit ticket}{110 -- 120 min}
Formulaire en ligne ou papier (annexe G) :
\begin{enumerate}
 \item #1
 \item #2
 \item Auto-évaluation : \emph{Je maîtrise / Je progresse / J'ai besoin d'aide.}
\end{enumerate}
Lire les réponses avant la séance suivante pour ajuster les groupes de besoin.
\end{etape}}

% Grille d'observation : #1 #2 #3 intitulés de colonnes
\newcommand{\grille}[3]{%
\subsection*{Grille d'observation (enseignant)}
\begin{tabularx}{\linewidth}{Y p{2.4cm} p{2.4cm} p{2.4cm} Y}
\toprule
Élève & #1 & #2 & #3 & Groupe suivant\\
\midrule
 & & & & \\ & & & & \\ & & & & \\ & & & & \\ & & & & \\
\bottomrule
\end{tabularx}}

% Signature de l'exit ticket
\newcommand{\autoeval}{Je me sens : \cb Je maîtrise \quad \cb Je progresse \quad \cb J'ai besoin d'aide}


\begin{document}
\titreblock{Les capteurs numériques et analogiques}{Semaine du 16 au 20 novembre 2026}

%==================================================================
\section{Informations générales}
\begin{tabularx}{\linewidth}{>{\bfseries}p{3.6cm}Y}
\toprule
Titre & Les capteurs numériques et analogiques\\
Classe / effectif & Grade 11, 11 à 20 élèves, niveaux hétérogènes, élèves à besoins particuliers\\
Durée & 2 heures (120 min)\\
Lieu & Salle informatique, un poste par élève\\
Matériel et logiciels &
\begin{itemize}
 \item Un poste par élève avec navigateur web.
 \item \textbf{Wokwi} (wokwi.com, gratuit) : projet ESP32 modèle avec un bouton poussoir (GPIO 4) et un potentiomètre (GPIO 34), moniteur série. Le code est fourni et les élèves modifient seulement des valeurs.
 \item Un tableur (Google Sheets ou LibreOffice Calc).
 \item Vidéoprojecteur ou tableau interactif.
\end{itemize}\\
Pré-requis &
Séance précédente : chaîne d'information (acquérir, traiter, communiquer). Notion de grandeur physique.\\
\bottomrule
\end{tabularx}

\begin{astuce}[Préparation avant le cours (environ 45 min)]
\begin{enumerate}
 \item Créer le projet Wokwi (ESP32, bouton sur GPIO 4, potentiomètre sur GPIO 34) et coller le code de l'annexe A. Tester le lien de partage.
 \item Préparer le tableur modèle (annexe B) pour relever les valeurs.
 \item Préparer les fichiers des annexes C à E et le formulaire de sortie.
\end{enumerate}
\end{astuce}

%==================================================================
\section{Objectifs d'apprentissage}
\begin{tcolorbox}[colback=gray!8,colframe=gray!60!black,arc=1mm,boxrule=0.5pt]
\textbf{Objectifs du programme}
\begin{itemize}
 \item Distinguer capteur analogique et numérique.
 \item Comprendre le principe d'acquisition.
\end{itemize}
\textbf{À la fin de la séance, l'élève est capable\,:}
\begin{itemize}
 \item de classer des capteurs courants en analogique ou numérique, avec une justification\,;
 \item d'expliquer qu'un signal analogique varie de façon continue alors qu'un signal numérique prend deux valeurs (0 ou 1)\,;
 \item de relever et d'interpréter les valeurs d'un capteur analogique et d'un capteur numérique dans un simulateur.
\end{itemize}
\textbf{Compétences numériques\,:} relever des mesures, utiliser un moniteur série, remplir un tableur.
\end{tcolorbox}

%==================================================================
\section{Déroulé minuté (120 min)}
\minutage{Interrupteur ou variateur de lumière ?}{Signal analogique, signal numérique, acquisition}{Bouton et potentiomètre en collectif, puis relevé de valeurs en binômes}{Trois niveaux de besoin}{Valoriser chaque groupe}{Trois questions sur poste}

%------------------------------------------------------------------
\begin{etape}{orange}{1. Accroche : interrupteur ou variateur ?}{0 -- 10 min}
\begin{enumerate}
 \item (5 min) Deux lampes : l'une avec un interrupteur (allumée ou éteinte), l'autre avec un variateur (toutes les intensités). Les élèves mettent en commun ce qu'on peut «~dire~» à chaque lampe.
 \item (5 min) Question-problème : \textbf{«~Comment un ordinateur lit-il un interrupteur ? Et un variateur ?~»}
\end{enumerate}
\end{etape}

%------------------------------------------------------------------
\begin{etape}{blue}{2. Mini-leçon : analogique et numérique}{10 -- 30 min}
\textbf{Support\,:} annexe A. Quatre temps de 5 min.

\textbf{Temps 1 : le signal numérique.} Deux états seulement : 0 (bas, environ 0~V) ou 1 (haut, environ 3,3~V sur ESP32). Exemple : bouton, détecteur de présence, fin de course.

\textbf{Temps 2 : le signal analogique.} Il prend une infinité de valeurs entre un minimum et un maximum. Exemple : potentiomètre, capteur de lumière, sonde de température.

\begin{center}
\begin{tabularx}{\linewidth}{>{\bfseries}p{3.4cm}Y Y}
\toprule
 & Numérique & Analogique\\
\midrule
Valeurs & 0 ou 1 & Toutes les valeurs d'une plage\\
Exemple de capteur & Bouton, détecteur de présence & Potentiomètre, thermistance\\
Lecture avec ESP32 & \texttt{digitalRead(broche)} & \texttt{analogRead(broche)}\\
Résultat & 0 ou 1 & 0 à 4095\\
\bottomrule
\end{tabularx}
\end{center}

\textbf{Temps 3 : l'acquisition.} Pour lire un signal analogique, le microcontrôleur utilise un convertisseur analogique-numérique (CAN) qui transforme la tension en un nombre. Sur ESP32, le CAN a 12 bits : la tension de 0 à 3,3~V devient une valeur de 0 à 4095.

\begin{center}
\begin{tikzpicture}[box/.style={draw,rounded corners,thick,minimum height=1cm,text width=2.5cm,align=center,font=\small},>={Stealth[length=3mm]},node distance=0.7cm]
\node[box,fill=orange!20] (a) {Grandeur physique};
\node[box,fill=blue!15,right=of a] (b) {Capteur\\ tension};
\node[box,fill=green!15,right=of b] (c) {CAN\\ nombre};
\node[box,fill=violet!15,right=of c] (d) {Programme};
\draw[->,thick] (a)--(b); \draw[->,thick] (b)--(c); \draw[->,thick] (c)--(d);
\end{tikzpicture}
\end{center}

\textbf{Temps 4 : comment choisir ?} Une information du type oui ou non demande un capteur numérique. Une mesure continue (température, distance, lumière) demande un capteur analogique.

\begin{astuce}
Garder les deux exemples (interrupteur, variateur) affichés. Faire manipuler physiquement un potentiomètre si possible. Lire les tableaux à voix haute.
\end{astuce}
\end{etape}

%------------------------------------------------------------------
\begin{etape}{green}{3. Pratique guidée}{30 -- 60 min}
\textbf{Fiche\,:} annexe B. Les élèves travaillent en binômes devant leur poste.

\textbf{3a. Bouton et potentiomètre en collectif (12 min).} Ouvrir le projet Wokwi. Appuyer sur le bouton et lire dans le moniteur série les valeurs 0 et 1. Tourner le potentiomètre et observer les valeurs entre 0 et 4095.

\textbf{3b. Relevé de valeurs en binômes (18 min).} Relever dans un tableur la valeur du potentiomètre pour cinq positions (0 \%, 25 \%, 50 \%, 75 \%, 100 \%) et calculer la tension correspondante : $U = \text{valeur} \times 3{,}3 \div 4095$.

\aides{Le tableur modèle avec une ligne déjà remplie}{Questions guides : \emph{Quelle est la valeur minimale ? Quelle est la valeur maximale ? La valeur augmente-t-elle de façon régulière ?}}{L'enseignant remplit une ligne avec le binôme, qui poursuit}

\begin{astuce}
Fournir le tableur avec la formule déjà posée pour les élèves qui ont des difficultés avec les calculs. Autoriser la calculatrice. Faire lire à voix haute les valeurs lues dans le moniteur.
\end{astuce}
\end{etape}

%------------------------------------------------------------------
\begin{etape}{violet}{4. Pratique autonome différenciée}{60 -- 95 min}
\groupes{Classer dix capteurs en analogique ou numérique avec une banque de mots (bouton, thermistance, détecteur de présence, potentiomètre, capteur de lumière, fin de course\ldots) et associer chacun à ce qu'il mesure. Pictogrammes en appui.}{Relever les valeurs du potentiomètre pour sept positions, calculer les tensions dans le tableur, tracer le graphique valeur en fonction de la position et interpréter (la relation est-elle linéaire ?).}{Comparer deux capteurs dans Wokwi (potentiomètre et capteur de lumière ou de température fourni), proposer un seuil pour un éclairage automatique, et expliquer pourquoi un capteur numérique de type «~nuit ou jour~» est parfois suffisant.}

\medskip
\textbf{Plan B si Wokwi est indisponible\,:} utiliser les valeurs relevées par l'enseignant (tableau imprimé) pour le tableur.
\end{etape}

%------------------------------------------------------------------
\miseencommun{un élève présente son classement et justifie deux choix. Valoriser l'argument.}{un binôme présente son graphique et explique la forme de la courbe. Valoriser l'interprétation.}{un élève présente son seuil et sa comparaison de capteurs. Valoriser l'esprit critique.}{\textit{un capteur numérique donne 0 ou 1, un capteur analogique donne une valeur qui varie ; le CAN convertit la tension en nombre.}}

\cloture{Un détecteur de présence est-il analogique ou numérique ? Pourquoi ?}{Que vaut la tension si \texttt{analogRead} renvoie 2048 sur un ESP32 (3,3~V, 12 bits) ?}

%==================================================================
% ANNEXES
%==================================================================
\newpage
\section*{Annexe A : Support résumé de la mini-leçon}
\entete{blue}{Capteurs numériques et analogiques}

\begin{tcolorbox}[colback=yellow!10,colframe=yellow!50!black]
\textbf{Numérique} : 0 ou 1. \textbf{Analogique} : toutes les valeurs d'une plage. Le \textbf{CAN} transforme une tension en nombre. ESP32 : 12 bits, de 0 à 4095 pour 0 à 3,3~V.
\end{tcolorbox}

\begin{tabularx}{\linewidth}{>{\bfseries}p{3.4cm}Y Y}
\toprule
 & Numérique & Analogique\\
\midrule
Valeurs & 0 ou 1 & 0 à 4095\\
Exemples & Bouton, présence & Potentiomètre, lumière\\
Instruction & \texttt{digitalRead} & \texttt{analogRead}\\
\bottomrule
\end{tabularx}

\begin{arduino}[Code fourni (Wokwi)]
void setup() {
  Serial.begin(115200);
  pinMode(4, INPUT_PULLUP);       // bouton
}

void loop() {
  int bouton = digitalRead(4);    // numerique : 0 ou 1
  int pot = analogRead(34);       // analogique : 0 a 4095
  Serial.print("bouton = ");
  Serial.print(bouton);
  Serial.print("   pot = ");
  Serial.println(pot);
  delay(300);
}
\end{arduino}

\textbf{Remarque :} avec \texttt{INPUT\_PULLUP}, le bouton relâché donne 1 et le bouton appuyé donne 0.

%==================================================================
\newpage
\section*{Annexe B : Fiche de pratique guidée}
\entete{green}{Pratique guidée : lire deux capteurs}

\subsection*{3a. Observations}
Bouton relâché : \dotfill\quad Bouton appuyé : \dotfill\\
Valeur minimale du potentiomètre : \dotfill\quad Valeur maximale : \dotfill

\subsection*{3b. Relevé de valeurs}
\begin{tabularx}{\linewidth}{>{\bfseries}p{3.4cm}Y Y}
\toprule
Position du potentiomètre & Valeur lue & Tension $U = \text{valeur} \times 3{,}3 \div 4095$ (V)\\
\midrule
0 \% & & \\ 25 \% & & \\ 50 \% & & \\ 75 \% & & \\ 100 \% & & \\
\bottomrule
\end{tabularx}

\begin{tcolorbox}[colback=yellow!10,colframe=yellow!50!black,title={Cartes d'aide à utiliser dans l'ordre},fonttitle=\bfseries\small]
\textbf{Aide 1 :} regarde le tableau de l'annexe A et la ligne déjà remplie du tableur.\\
\textbf{Aide 2 :} Quelle est la valeur minimale ? Quelle est la valeur maximale ? La valeur augmente-t-elle de façon régulière ?\\
\textbf{Aide 3 :} demande à l'enseignant de remplir une ligne avec toi.
\end{tcolorbox}

%==================================================================
\newpage
\section*{Annexe C : Niveau A (en dessous du niveau)}
\entete{orange}{Niveau A : numérique ou analogique ?}

\textbf{Banque :} bouton, thermistance, détecteur de présence, potentiomètre, capteur de lumière, fin de course.

\begin{tabularx}{\linewidth}{>{\bfseries}p{4cm}c c Y}
\toprule
Capteur & Numérique & Analogique & Ce qu'il mesure\\
\midrule
Bouton & \cb & \cb & \\
Thermistance & \cb & \cb & \\
Détecteur de présence & \cb & \cb & \\
Potentiomètre & \cb & \cb & \\
Capteur de lumière & \cb & \cb & \\
Fin de course & \cb & \cb & \\
\bottomrule
\end{tabularx}

\subsection*{Vrai ou faux ?}
1. Un signal numérique prend deux valeurs. \cb Vrai \cb Faux\\
2. \texttt{analogRead} renvoie 0 ou 1. \cb Vrai \cb Faux

%==================================================================
\newpage
\section*{Annexe D : Niveau B (au niveau)}
\entete{blue}{Niveau B : relevé et graphique}

\begin{tabularx}{\linewidth}{>{\bfseries}p{3cm}Y Y}
\toprule
Position & Valeur lue & Tension (V)\\
\midrule
0 \% & & \\ 20 \% & & \\ 40 \% & & \\ 50 \% & & \\ 60 \% & & \\ 80 \% & & \\ 100 \% & & \\
\bottomrule
\end{tabularx}

Trace dans le tableur un graphique en nuage de points (position en abscisse, valeur en ordonnée).

\subsection*{Questions}
\begin{enumerate}
 \item La relation est-elle proche d'une droite ? \dotfill
 \item Quelle tension correspond à la valeur 1024 ? \dotfill
\end{enumerate}

%==================================================================
\newpage
\section*{Annexe E : Niveau C (au-dessus du niveau)}
\entete{violet}{Niveau C : choisir le bon capteur}

\textbf{Situation :} un éclairage de couloir doit s'allumer la nuit et seulement si quelqu'un passe.

\begin{tabularx}{\linewidth}{>{\bfseries}p{3.4cm}Y Y Y}
\toprule
Besoin & Capteur choisi & Numérique ou analogique ? & Justification\\
\midrule
Détecter une personne & & & \\
Savoir s'il fait nuit & & & \\
\bottomrule
\end{tabularx}

\begin{enumerate}
 \item Propose un seuil pour la valeur du capteur de lumière (0 à 4095) et justifie : \lig
 \item Pourquoi un capteur numérique «~nuit ou jour~» peut-il suffire ? Quel avantage et quelle limite ? \lig \lig
\end{enumerate}

\textbf{Défi bonus :} explique ce que l'on perd quand on convertit une mesure analogique en 12 bits (résolution).

%==================================================================
\newpage
\section*{Annexe F : Corrections et attendus (enseignant)}

\subsection*{Annexe B}
Bouton relâché : 1 ; appuyé : 0 (INPUT\_PULLUP). Potentiomètre : minimum 0, maximum 4095. Tensions : 0 ; environ 0,8 ; 1,65 ; 2,5 ; 3,3 V.

\subsection*{Annexe C}
Numériques : bouton, détecteur de présence, fin de course. Analogiques : thermistance (température), potentiomètre (position), capteur de lumière (luminosité). Vrai, Faux.

\subsection*{Annexe D}
Relation à peu près linéaire. Tension pour 1024 : $1024 \times 3{,}3 \div 4095 \approx 0{,}83$~V.

\subsection*{Annexe E}
Présence : capteur infrarouge passif (numérique). Nuit : capteur de lumière (analogique avec seuil, par exemple 1000) ou capteur crépusculaire (numérique). Un capteur numérique est plus simple mais ne permet pas de régler finement le seuil. Résolution : 4096 niveaux, soit environ 0,8 mV par pas.\\
\textbf{Critères de réussite :} classement juste, calcul de tension correct, seuil justifié.

%==================================================================
\newpage
\section*{Annexe G : Exit ticket et grille d'observation}
\entete{gray}{Exit ticket : capteurs}

\begin{enumerate}
 \item Un détecteur de présence est-il analogique ou numérique ? Pourquoi ?\\ \lig
 \item Que vaut la tension si \texttt{analogRead} renvoie 2048 sur un ESP32 (3,3~V, 12 bits) ?\\ \lig
 \item \autoeval
\end{enumerate}
\textbf{Attendus :} 1. Numérique : il donne deux états (présence ou non). 2. Environ 1,65~V (la moitié de 3,3~V).

\bigskip
\grille{Classement correct}{Tension calculée}{Interprétation}

\end{document}
